\documentclass[12pt]{article}
\usepackage[utf8]{inputenc}
\usepackage[french]{babel}
\usepackage{geometry}
\usepackage{enumitem}
\usepackage{titlesec}
\usepackage{hyperref}
\usepackage{xcolor}
\usepackage{tikz}
\usetikzlibrary{shapes.geometric, arrows.meta, positioning, calc}

\usepackage{listings} \lstset{ language=SQL, basicstyle=\ttfamily\small, keywordstyle=\color{blue}\bfseries, commentstyle=\color{gray}, stringstyle=\color{red}, frame=single, breaklines=true }
\geometry{a4paper, margin=2.5cm}
\titleformat{\section}{\large\bfseries\color{blue}}{\thesection}{1em}{}
\titleformat{\subsection}{\normalsize\bfseries\color{teal}}{\thesubsection}{1em}{}
\author{
  Talel Laarif
  \and
Kallel Mohamed Amine
  \and
  Djemmali Yassine
}
\title{\textbf{Modélisation du Système de Pharmacie}}


\begin{document}

\maketitle

\section*{Introduction}
Ce document présente la modélisation du système de gestion d'une pharmacie selon deux niveaux :
\begin{itemize}
  \item Le \textbf{Modèle Conceptuel des Données (MCD)} : basé sur les entités et associations.
  \item Le \textbf{Modèle Logique des Données (MLD)} : basé sur les relations et les clés.
\end{itemize}

\section{Modèle Conceptuel des Données (MCD)}

\subsection*{Entités}
\begin{itemize}
  \item \textbf{Produit} : idProduit (AUTO\_INCREMENT), nomProduit, marque, quantite, quantiteMinimale, prix, TVA, type
  \item \textbf{Fournisseur} : idFournisseur, nom, prenom, numeroTelephone, adresseEmail
  \item \textbf{Commande} : idCommande, dateCommande, periodeReception, recu, dateReception
  \item \textbf{Client} : idClient, nom, adresse, email, nTelephone
  \item \textbf{Vente} : idVente, dateFacture
  \item \textbf{Employe} : idEmploye, nom, prenom, login, motDePasse, email, dateEmbauche, actif
\end{itemize}

\subsection*{Associations}
\begin{itemize}
  \item \textbf{Produit-Fournisseur} (n,n) : lie plusieurs produits à plusieurs fournisseurs
  \item \textbf{LigneCommande} (n,n) : chaque commande contient plusieurs produits
  \item \textbf{LigneVente} (n,n) : chaque vente contient plusieurs produits
  \item \textbf{Vente-Client} (n,1) : chaque vente est liée à un seul client
  \item \textbf{Commande-Fournisseur} (n,1) : chaque commande est passée à un seul fournisseur

\end{itemize}


\subsection*{Diagramme MCD}
\begin{center}
\begin{tikzpicture}[
  entity/.style={rectangle, draw=blue!70, thick, minimum width=3.8cm, minimum height=1.2cm, align=center},
  relation/.style={ellipse, draw=red!70, thick, minimum width=2.2cm, minimum height=1cm, align=center},
  arrow/.style={-{Stealth[length=3mm]}, thick},
  node distance=2cm and 5cm
]

% Entities
\node[entity] (produit) {Produit\\ \scriptsize \textbf{idProduit} (AUTO\_INC), nomProduit, marque, quantite, 
\newline, prix, TVA, type};
\node[entity, right=2cm of produit] (fournisseur) {Fournisseur\\ \scriptsize \textbf{idFournisseur}, nom, prenom, tel, email};
\node[entity, below=5cm of produit] (commande) {Commande\\ \scriptsize \textbf{idCommande}, date, recu};
\node[entity, below=3cm of fournisseur] (client) {Client\\ \scriptsize \textbf{idClient}, nom, adresse, email, tel};
\node[entity, below=3cm of client] (vente) {Vente\\ \scriptsize \textbf{idVente}, dateFacture};
\node[entity, below=2cm of commande] (employe) {Employe\\ \scriptsize \textbf{idEmploye}, login, nom, prenom, actif};

% Relations
% Example: place relation nodes 1.2cm above or below midpoint
\node[relation] (pf) at ($(produit)!0.5!(fournisseur)+(0.2,-1.7)$) {Produit-Fournisseur};
\node[relation] (lc) at ($(produit)!0.5!(commande)+(0,0)$) {LigneCommande};
\node[relation] (lv) at ($(produit)!0.5!(vente)+(0,-1.2)$) {LigneVente};
\node[relation] (vc) at ($(vente)!0.5!(client)+(0,0)$) {Vente-Client};
\node[relation] (cf) at ($(commande)!0.5!(fournisseur)+(0,-0.5)$) {Commande-Fournisseur};




% Arrows with cardinalities
\draw[arrow] (produit) -- node[above] {n} (pf);
\draw[arrow] (fournisseur) -- node[above] {n} (pf);

\draw[arrow] (produit) -- node[left] {n} (lc);
\draw[arrow] (commande) -- node[right] {n} (lc);
\draw[arrow, shorten >=2pt, shorten <=2pt] (commande) -- (lc);


\draw[arrow] (produit) -- node[left] {n} (lv);
\draw[arrow] (vente) -- node[right] {n} (lv);

\draw[arrow] (vente) -- node[right] {n} (vc);
\draw[arrow] (client) -- node[left] {1} (vc);
\draw[arrow, shorten >=2pt, shorten <=2pt] (commande) -- node[below] {n} (cf);
\draw[arrow, shorten >=2pt, shorten <=2pt] (fournisseur) -- node[below] {1} (cf);

\end{tikzpicture}

\end{center}

\section{Modèle Logique des Données (MLD)}

\begin{itemize}[label=--]
  \item \textbf{Produit}(\underline{idProduit}, nomProduit, marque, quantite, quantiteMinimale, prix, TVA, type)\\
    \textit{Contrainte : idProduit AUTO\_INCREMENT, quantiteMinimale $\geq$ 0}
  \item \textbf{Fournisseur}(\underline{idFournisseur}, nom, prenom, numeroTelephone,adresseEmail)
  \item \textbf{Produit\_Fournisseur}(\underline{idProduit}, \underline{idFournisseur})\\
    \textit{Clés étrangères : idProduit → Produit, idFournisseur → Fournisseur}
  \item \textbf{Commande}(\underline{idCommande}, \underline{idFournisseur}, dateCommande, periodeReception, recu, dateReception)\\
    \textit{Clé étrangère : idFournisseur → Fournisseur}
  \item \textbf{LigneCommande}(\underline{idCommande}, \underline{idProduit}, quantite, prixAchat)\\
    \textit{Clés étrangères : idCommande → Commande, idProduit → Produit}
  \item \textbf{Vente}(\underline{idVente}, dateFacture,\underline{idClient},)
  \item \textbf{LigneVente}(\underline{idVente}, \underline{idProduit}, quantite, prixUnite)\\
    \textit{Clés étrangères : idVente → Facture, idProduit → Produit}
 \item \textbf{Client}(\underline{idClient}, nom,adresse,email,nTelephone)
 \item \textbf{Employe}(\underline{idEmploye}, nom,prenom,email,login,mot de passe,dateEmbauche,actif)
\end{itemize}
\section{Initialisation SQL de la Base de Données}

\begin{lstlisting}



CREATE TABLE Produit (
  idProduit INT NOT NULL AUTO_INCREMENT PRIMARY KEY,
  nomProduit VARCHAR(100),
  marque VARCHAR(50),
  quantite INT,
  quantiteMinimale INT DEFAULT 10 NOT NULL,
  prix DECIMAL(10,2),
  TVA DECIMAL(5,2),
  type VARCHAR(50)
);

CREATE TABLE Fournisseur (
  idFournisseur INT NOT NULL AUTO_INCREMENT PRIMARY KEY,
  nom VARCHAR(50),
  prenom VARCHAR(50),
  numeroTelephone VARCHAR(20),
  adresseEmail VARCHAR(100)
);

CREATE TABLE Produit_Fournisseur (
  idProduit INT,
  idFournisseur INT,
  PRIMARY KEY (idProduit, idFournisseur),
  FOREIGN KEY (idProduit) REFERENCES Produit(idProduit),
  FOREIGN KEY (idFournisseur) REFERENCES Fournisseur(idFournisseur)
);

CREATE TABLE Commande (
  idCommande INT NOT NULL AUTO_INCREMENT PRIMARY KEY,
  idFournisseur INT,
  dateCommande DATE,
  periodeReception INT,
  recu BOOLEAN DEFAULT FALSE,
  dateReception DATE,
  FOREIGN KEY (idFournisseur) REFERENCES Fournisseur(idFournisseur)
);

CREATE TABLE LigneCommande (
  idCommande INT,
  idProduit INT,
  quantite INT,
  prixAchat DECIMAL(10,2),
  PRIMARY KEY (idCommande, idProduit),
  FOREIGN KEY (idCommande) REFERENCES Commande(idCommande),
  FOREIGN KEY (idProduit) REFERENCES Produit(idProduit)
);

CREATE TABLE Client (
  idClient INT NOT NULL AUTO_INCREMENT PRIMARY KEY,
  nom VARCHAR(50),
  adresse VARCHAR(100),
  email VARCHAR(100),
  nTelephone VARCHAR(20)
);

CREATE TABLE Vente (
  idVente INT NOT NULL AUTO_INCREMENT PRIMARY KEY,
  dateFacture DATE,
  idClient INT,
  FOREIGN KEY (idClient) REFERENCES Client(idClient)
);

CREATE TABLE LigneVente (
  idVente INT,
  idProduit INT,
  quantite INT,
  prixUnite DECIMAL(10,2),
  PRIMARY KEY (idVente, idProduit),
  FOREIGN KEY (idVente) REFERENCES Vente(idVente),
  FOREIGN KEY (idProduit) REFERENCES Produit(idProduit)
);

CREATE TABLE Employe (
  idEmploye INT PRIMARY KEY AUTO_INCREMENT,
  nom VARCHAR(50) NOT NULL,
  prenom VARCHAR(50) NOT NULL,
  login VARCHAR(50) UNIQUE NOT NULL,
  motDePasse VARCHAR(255) NOT NULL,
  email VARCHAR(100),
  dateEmbauche DATE,
  actif BOOLEAN DEFAULT TRUE
);

CREATE INDEX idx_login ON Employe(login);

INSERT INTO Employe (nom, prenom, login, motDePasse, email, dateEmbauche, actif) 
VALUES ('Admin', 'Pharmacie', 'admin', 'admin123', 'admin@pharmacie.com', NOW(), TRUE);

\end{lstlisting}



\section{Diagramme UML - Couche Modèle (Objets Métier)}
\begin{tikzpicture}[
  class/.style={rectangle, draw=blue!70, thick, fill=blue!5, minimum width=4.5cm, minimum height=3cm, align=left, font=\scriptsize},
  relation/.style={-Stealth, thick},
  node distance=5.5cm
]

% Row 1
\node[class] (produit) {
  \textbf{Produit}\\
  \rule{4.3cm}{0.4pt}\\
  - idProduit: int\\
  - nomProduit: String\\
  - marque: String\\
  - quantite: int\\
  - quantiteMinimale: int\\
  - prix: double\\
  - tva: double\\
  - type: String\\
  \rule{4.3cm}{0.4pt}\\
  + estEnStockCritique(): boolean\\
  + augmenterQuantite(q: int): void\\
  + diminuerQuantite(q: int): void
};

\node[class, right of=produit] (fournisseur) {
  \textbf{Fournisseur}\\
  \rule{4.3cm}{0.4pt}\\
  - idFournisseur: int\\
  - nom: String\\
  - prenom: String\\
  - numeroTelephone: String\\
  - adresseEmail: String\\
  \rule{4.3cm}{0.4pt}\\
  + getId(): int\\
  + getNom(): String\\
  + setNom(nom: String): void\\
  + getEmail(): String
};

% Row 2
\node[class, below of=produit] (commande) {
  \textbf{Commande}\\
  \rule{4.3cm}{0.4pt}\\
  - idCommande: int\\
  - idFournisseur: int\\
  - dateCommande: LocalDate\\
  - periodeReception: int\\
  - recu: boolean\\
  - dateReception: LocalDate\\
  \rule{4.3cm}{0.4pt}\\
  + getId(): int\\
  + marquerRecue(): void\\
  + estEnRetard(): boolean
};

\node[class, right of=commande] (lignecmd) {
  \textbf{LigneCommande}\\
  \rule{4.3cm}{0.4pt}\\
  - idCommande: int\\
  - idProduit: int\\
  - quantite: int\\
  - prixAchat: double\\
  \rule{4.3cm}{0.4pt}\\
  + getQuantite(): int\\
  + getPrixAchat(): double\\
  + getSousTotal(): double
};

% Row 3
\node[class, below of=commande] (client) {
  \textbf{Client}\\
  \rule{4.3cm}{0.4pt}\\
  - idClient: int\\
  - nom: String\\
  - adresse: String\\
  - email: String\\
  - nTelephone: String\\
  \rule{4.3cm}{0.4pt}\\
  + getId(): int\\
  + getNom(): String\\
  + setNom(nom: String): void
};

\node[class, right of=client] (vente) {
  \textbf{Vente}\\
  \rule{4.3cm}{0.4pt}\\
  - idVente: int\\
  - dateFacture: LocalDate\\
  - idClient: int\\
  \rule{4.3cm}{0.4pt}\\
  + getId(): int\\
  + getIdClient(): int\\
  + getDateFacture(): LocalDate\\
  + calculerTotal(): double
};

% Row 4
\node[class, below of=client] (lignevte) {
  \textbf{LigneVente}\\
  \rule{4.3cm}{0.4pt}\\
  - idVente: int\\
  - idProduit: int\\
  - quantite: int\\
  - prixUnite: double\\
  \rule{4.3cm}{0.4pt}\\
  + getQuantite(): int\\
  + getPrixUnite(): double\\
  + getSousTotal(): double
};

\node[class, right of=lignevte] (employe) {
  \textbf{Employe}\\
  \rule{4.3cm}{0.4pt}\\
  - idEmploye: int\\
  - nom: String\\
  - prenom: String\\
  - login: String\\
  - motDePasse: String\\
  - email: String\\
  - dateEmbauche: LocalDate\\
  - actif: boolean\\
  \rule{4.3cm}{0.4pt}\\
  + getId(): int\\
  + getLogin(): String\\
  + estActif(): boolean
};

% Relationships
\draw[relation] (commande) -- node[above, sloped, font=\tiny] {utilise} (produit);
\draw[relation] (commande) -- node[above, font=\tiny] {passe à} (fournisseur);
\draw[relation] (commande) -- node[right, font=\tiny] {contient} (lignecmd);
\draw[relation] (vente) -- node[above, font=\tiny] {pour} (client);
\draw[relation] (vente) -- node[right, font=\tiny] {contient} (lignevte);

\end{tikzpicture}
\end{center}

\section{Diagramme UML - Couche DAO (Accès aux Données)}

\begin{center}
\begin{tikzpicture}[
  classdao/.style={
    rectangle, draw=blue!70, thick, fill=blue!5,
    minimum width=4.8cm, minimum height=3cm,
    align=left, font=\scriptsize
  },
  relation/.style={-Stealth, thick},
  node distance=5.5cm
]

% ===== Row 1 =====
\node[classdao] (produitdao) {
  \textbf{ProduitDAO}\\
  \rule{4.6cm}{0.4pt}\\
  + save(p:Produit): void\\
  + update(p:Produit): void\\
  + findById(id:int): Produit\\
  + findAll(): List<Produit>\\
  + delete(id:int): void
};

\node[classdao, right of=produitdao] (fournisseurdao) {
  \textbf{FournisseurDAO}\\
  \rule{4.6cm}{0.4pt}\\
  + save(f:Fournisseur): void\\
  + update(f:Fournisseur): void\\
  + findById(id:int): Fournisseur\\
  + findAll(): List<Fournisseur>\\
  + delete(id:int): void
};

% ===== Row 2 =====
\node[classdao, below of=produitdao] (commandedao) {
  \textbf{CommandeDAO}\\
  \rule{4.6cm}{0.4pt}\\
  + save(c:Commande): int\\
  + update(c:Commande): void\\
  + findById(id:int): Commande\\
  + findAll(): List<Commande>\\
  + delete(id:int): void
};

\node[classdao, right of=commandedao] (lignecommandedao) {
  \textbf{LigneCommandeDAO}\\
  \rule{4.6cm}{0.4pt}\\
  + save(lc:LigneCommande): void\\
  + findByCommande(id:int): List<LigneCommande>\\
  + deleteByCommande(id:int): void
};

% ===== Row 3 =====
\node[classdao, below of=commandedao] (clientdao) {
  \textbf{ClientDAO}\\
  \rule{4.6cm}{0.4pt}\\
  + save(c:Client): void\\
  + update(c:Client): void\\
  + findById(id:int): Client\\
  + findAll(): List<Client>\\
  + delete(id:int): void
};

\node[classdao, right of=clientdao] (ventedao) {
  \textbf{VenteDAO}\\
  \rule{4.6cm}{0.4pt}\\
  + save(v:Vente): int\\
  + findById(id:int): Vente\\
  + findAll(): List<Vente>\\
  + delete(id:int): void
};

% ===== Row 4 =====
\node[classdao, below of=clientdao] (ligneventedao) {
  \textbf{LigneVenteDAO}\\
  \rule{4.6cm}{0.4pt}\\
  + save(lv:LigneVente): void\\
  + findByVente(id:int): List<LigneVente>\\
  + deleteByVente(id:int): void
};

\node[classdao, right of=ligneventedao] (employedao) {
  \textbf{EmployeDAO}\\
  \rule{4.6cm}{0.4pt}\\
  + save(e:Employe): void\\
  + findByLogin(login:String): Employe\\
  + findAll(): List<Employe>\\
  + deactivate(id:int): void
};

% ===== Relations (clean & minimal) =====
\draw[relation] (produitdao) -- (fournisseurdao);
\draw[relation] (commandedao) -- (lignecommandedao);
\draw[relation] (ventedao) -- (ligneventedao);
\draw[relation] (ventedao) -- (clientdao);

\end{tikzpicture}
\end{center}



\end{document}